% ---------- 2. History ----------
\begin{frame}{A short history of AI, 1943 to 2026}
\begin{sfig}[0.78]
\begin{tikzpicture}[x=1mm,y=1mm]
  \draw[very thick,tdMuted] (0,0) -- (152,0);
  \foreach \y in {1950,1960,1970,1980,1990,2000,2010,2020}{\pgfmathsetmacro{\xx}{\y<2010 ? (\y-1940)*1.1 : 77+(\y-2010)*4.4}\draw[tdMuted] (\xx,-1.2) -- (\xx,1.2); \node[font=\tiny,text=tdMuted] at (\xx,-3.5) {\y};}
  \fill[tdRed!18] ({(1974-1940)*1.1},-1.4) rectangle ({(1980-1940)*1.1},1.4);
  \fill[tdRed!18] ({(1987-1940)*1.1},-1.4) rectangle ({(1995-1940)*1.1},1.4);
  \node[font=\tiny,text=tdRed] at ({(1977-1940)*1.1},-7) {winter};
  \node[font=\tiny,text=tdRed] at ({(1991-1940)*1.1},-7) {winter};
  \foreach \y/\t/\c/\h in {1943/{McCulloch--Pitts\\ neuron}/tdBlue/10,1950/{Turing:\\ imitation game}/tdAmber/24,1956/{Dartmouth\\ workshop}/tdGreen/10,1958/{Perceptron}/tdBlue/-14,1966/{ELIZA}/tdAmber/14,1997/{Deep Blue\\ beats Kasparov}/tdGreen/14,2012/{AlexNet,\\ deep learning}/tdBlue/12,2016/{AlphaGo}/tdGreen/-12,2017/{Transformer}/tdPurple/28,2020/{GPT-3}/tdPurple/-26,2022/{ChatGPT}/tdRed/12,2024/{reasoning\\ models}/tdAmber/-12,2025/{agents}/tdBlue/-26}{
    \pgfmathsetmacro{\xx}{\y<2010 ? (\y-1940)*1.1 : 77+(\y-2010)*4.4}
    \fill[\c] (\xx,0) circle (1.1mm);
    \ifnum\h<0 \draw[\c!60] (\xx,0) -- (\xx,{\h+5}); \node[font=\scriptsize,align=center,text=\c!70!black,anchor=north] at (\xx,{\h+5}) {\textbf{\y}\\ \t};
    \else \draw[\c!60] (\xx,0) -- (\xx,{\h-1}); \node[font=\scriptsize,align=center,text=\c!70!black,anchor=south] at (\xx,{\h-1}) {\textbf{\y}\\ \t};\fi
  }
\end{tikzpicture}
\end{sfig}
\vspace{-1mm}
{\scriptsize Horizontal scale compressed before 2010 and stretched after. Sources: McCulloch and Pitts 1943; Turing 1950; Dartmouth 1956; Rosenblatt 1958; Weizenbaum 1966; Deep Blue 1997; AlexNet 2012; AlphaGo 2016; Vaswani et al.\ 2017; GPT-3 2020; ChatGPT 30 Nov 2022. Reasoning models (2024) and agents (2025): see the ``AI now'' part. Winter dates approximate.}
\end{frame}

\begin{frame}{Dartmouth, summer 1956: the field gets its name}
\begin{columns}[T]
\begin{column}{0.52\textwidth}
\small
John McCarthy, Marvin Minsky, Nathaniel Rochester and Claude Shannon proposed a two-month study at Dartmouth College. The proposal introduced the phrase \textbf{``artificial intelligence''}.

The conjecture at its centre:
\vspace{1mm}

\qcard{tdBlue}{\linewidth-10pt}{``every aspect of learning or any other feature of intelligence can in principle be so precisely described that a machine can be made to simulate it.''}
\vspace{1mm}

They expected a ``significant advance'' from a small group over one summer.
\end{column}
\begin{column}{0.45\textwidth}
\begin{sfig}[0.46]\begin{tikzpicture}
  \node[card=tdBlue,text width=3.0cm,minimum height=1.8cm] (a) at (0,1.2) {\textbf{Symbolic AI}\\{\footnotesize logic, search, expert systems}};
  \node[card=tdGreen,text width=3.0cm,minimum height=1.8cm] (b) at (0,-1.2) {\textbf{Connectionism}\\{\footnotesize neural networks that learn}};
  \node[card=tdRed,text width=2.6cm,minimum height=1.2cm] (c) at (4.4,1.2) {\footnotesize two winters};
  \node[card=tdPurple,text width=2.6cm,minimum height=1.2cm] (d) at (4.4,-1.2) {\footnotesize deep learning and LLMs};
  \draw[arr] (a) -- (c); \draw[arr] (b) -- (d);
\end{tikzpicture}\end{sfig}
{\small The two traditions took turns being fashionable, and the current systems are connectionist at heart.}
\end{column}
\end{columns}
\src{McCarthy, Minsky, Rochester, Shannon, ``A Proposal for the Dartmouth Summer Research Project on Artificial Intelligence'', 31 Aug 1955.}
\end{frame}

\begin{frame}{Hype, winters and comebacks}
\begin{columns}[T]
\begin{column}{0.56\textwidth}
\begin{sfig}[0.62]
\begin{tikzpicture}
  \draw[->] (0,0) -- (8.6,0) node[right,font=\scriptsize] {time};
  \draw[->] (0,0) -- (0,3.4) node[above,font=\scriptsize] {expectations};
  \draw[very thick,tdBlue,smooth] plot coordinates {(0.2,0.4) (0.9,2.2) (1.6,3.0) (2.4,1.2) (3.0,0.6) (3.6,1.8) (4.3,2.7) (5.1,0.7) (6.0,0.5) (6.8,1.5) (7.4,2.6) (8.2,3.2)};
  \node[font=\tiny,align=center,anchor=south] at (1.6,3.0) {1950s--60s\\ first boom};
  \node[font=\tiny,align=center,text=tdRed,anchor=north] at (3.0,0.55) {first\\ winter};
  \node[font=\tiny,align=center,anchor=south] at (4.3,2.7) {1980s expert\\ systems};
  \node[font=\tiny,align=center,text=tdRed,anchor=north] at (5.8,0.45) {second\\ winter};
  \node[font=\tiny,align=center,anchor=south] at (8.2,3.2) {deep learning\\ and LLMs};
\end{tikzpicture}
\end{sfig}
{\scriptsize\illus}
\end{column}
\begin{column}{0.41\textwidth}
\small
\textbf{Why winters happened}
\begin{itemize}\setlength\itemsep{2pt}
\item promises ran ahead of results;
\item early methods did not scale (combinatorial explosion, brittle rules);
\item funders cut budgets, notably after the 1973 Lighthill report in the UK.
\end{itemize}
\textbf{What ended them:} more data, more compute, better learning methods.
\end{column}
\end{columns}
\discuss{We are in a boom. Which of the winter warning signs would you watch for today?}
\src{Lighthill, ``Artificial Intelligence: A General Survey'' (1973); Russell and Norvig, AIMA, 4th ed., ch.\ 1 (history).}
\end{frame}

\begin{frame}{The deep-learning turn}
\begin{columns}[T]
\begin{column}{0.5\textwidth}
\small
\begin{itemize}\setlength\itemsep{4pt}
\item[\key{2012}] \textbf{AlexNet} wins the ImageNet image-recognition contest by a wide margin, using GPUs and a deep convolutional network.
\item[\key{2016}] \textbf{AlphaGo} beats Lee Sedol 4--1 at Go, a game once thought decades away.
\item[\key{2017}] The \textbf{Transformer} (``Attention Is All You Need'') replaces recurrence with attention and scales well.
\item[\key{2020}] \textbf{GPT-3}: 175 billion parameters, few-shot learning from a prompt.
\item[\key{2022}] \textbf{ChatGPT} (30 Nov): a chat interface brings LLMs to hundreds of millions of people.
\end{itemize}
\end{column}
\begin{column}{0.47\textwidth}
\begin{sfig}[0.42]
\begin{tikzpicture}
\begin{axis}[width=7.4cm,height=4.4cm,ymode=log,xmin=2011,xmax=2021,ymin=2e7,ymax=1e12,xlabel={year},ylabel={parameters},xtick={2012,2017,2019,2020},xticklabels={2012,2017,2019,2020},ymajorgrids,grid style={tdLine}]
\addplot[only marks,mark=*,mark size=2.4pt,tdBlue] coordinates {(2012,6e7) (2017,6.5e7) (2019,1.5e9) (2020,1.75e11)};
\node[font=\tiny,anchor=south west] at (axis cs:2012.1,8e7) {AlexNet 60M};
\node[font=\tiny,anchor=north west] at (axis cs:2017.1,5e7) {Transformer (base) 65M};
\node[font=\tiny,anchor=south east] at (axis cs:2018.9,2.2e9) {GPT-2 1.5B};
\node[font=\tiny,anchor=north east] at (axis cs:2019.9,1.3e11) {GPT-3 175B};
\end{axis}
\end{tikzpicture}
\end{sfig}
{\scriptsize Parameter counts as reported in the original papers; recipe: big Transformer, next-token prediction, human feedback, scale.}
{\small The lesson many researchers took: \textbf{general methods that scale} beat hand-written knowledge (Sutton's ``bitter lesson'', 2019).}
\end{column}
\end{columns}
\src{Krizhevsky et al., NeurIPS 2012; Silver et al., Nature 2016; Vaswani et al., NeurIPS 2017; Brown et al., NeurIPS 2020; OpenAI, ChatGPT launch 30 Nov 2022; Sutton, ``The Bitter Lesson'' (2019).}
\end{frame}

\begin{frame}{From chatbots to reasoning models to agents}
\begin{sfig}[0.45]
\begin{tikzpicture}
  \node[card=tdBlue,minimum width=3.4cm,minimum height=2cm] (a) at (0,0) {\textbf{Pretrained LLM}\\{\scriptsize predicts next token}};
  \node[card=tdGreen,minimum width=3.4cm,minimum height=2cm] (b) at (4.3,0) {\textbf{Chat assistant}\\\shortstack{\scriptsize instructions +\\ human feedback}};
  \node[card=tdAmber,minimum width=3.4cm,minimum height=2cm] (c) at (8.6,0) {\textbf{Reasoning model}\\\shortstack{\scriptsize ``thinks'' longer at\\ test time}};
  \node[card=tdRed,minimum width=3.4cm,minimum height=2cm] (d) at (12.9,0) {\textbf{Agent}\\\shortstack{\scriptsize tools, code, browser,\\ long tasks}};
  \draw[arr] (a)--(b); \draw[arr] (b)--(c); \draw[arr] (c)--(d);
  \node[lbl] at (0,-1.6) {2018--2021}; \node[lbl] at (4.3,-1.6) {2022}; \node[lbl] at (8.6,-1.6) {2024--2025}; \node[lbl] at (12.9,-1.6) {2025--2026};
\end{tikzpicture}
\end{sfig}
\vspace{1mm}
\small
The ARC Prize 2025 analysis calls \textbf{refinement loops} (generate a candidate, check it against feedback, try again) ``the central theme'' of 2025 progress: the same idea behind reasoning models and agents, spending compute at \emph{inference} rather than only at training.
\discuss{If a model improves by ``thinking longer'', is that more intelligence, or more compute?}
\src{ARC Prize Foundation, ``ARC Prize 2025 Results and Analysis'', arcprize.org/blog/arc-prize-2025-results-analysis (Dec 2025); date ranges for the stages are approximate.}
\end{frame}
